% ============================================================================
% Case study: unconstrained output path in excel-reader-mcp::export_to_json.
% Every fact here is traceable:
%   provenance  ArchimedesCrypto/excel-reader-mcp @ 2c41e8a9, Set A, 19 stars
%   selection   eval/state/quarantine-bug20-20260911/gpt41-full/
%                 ArchimedesCrypto__excel-reader-mcp__export_to_json.log  (LLM_SELECTION_OK)
%   seed        eval/corpora_gpt41-full/.../export_to_json/seed_4.json    (LLM_SEEDGEN_OK)
%   sink        eval/targets/.../src/tools/export/export-to-json.ts:89-90
%   oracle      src/fuzzer.py:980  ("Silent Bug Alert")
%   overwrite   directed re-validation under the correct cwd, 2026-09-11
% See nazmul.todo, "RE-TRIAGE ... V1", for the full evidence chain.
% ============================================================================

\subsection{Case study: an unconstrained output path that no crash oracle can see}
\label{sec:casestudy-excel}

This example is representative of \project{}'s contribution precisely because
nothing about it looks wrong to a conventional fuzzer. The offending argument
is a syntactically valid string, it passes the tool's own type check, it is
never passed to a shell, and the tool returns a normal success response with no
exception, no crash, and no error payload. A malformed-input or crash-based
oracle observes a clean, successful call. The violation surfaces only when the
string is interpreted as a filesystem path and compared against what the tool's
documented interface promises---exactly the LLM-to-tool semantic gap
\project{} is built to expose.

\para{Documented interface.}
The tool \texttt{export\_to\_json} from \texttt{excel-reader-mcp} is described
in its metadata as ``Export a sheet to JSON format.'' Its \texttt{outputPath}
argument is documented as a string: ``Path to save JSON file (optional, returns
content if not provided).'' The documented contract is therefore that the tool
serialises \emph{its own} spreadsheet output; \texttt{outputPath} names where
that output is saved. Nothing in the interface indicates that a caller may
direct the write outside a working directory or onto a pre-existing file.

\para{What \project{} receives and selects ($A$).}
\project{} passes this metadata---name, description, and the argument
schema---to the sanitizer-selection stage. From the corpus, the LLM selects six
file-oriented sanitizers, driven by the two path-bearing arguments
(\texttt{filePath}, \texttt{outputPath}); the applicable check here is
\texttt{check\_path\_traversal}, whose corpus description is ``Used for tools
that manipulate, read, or write files on the file system.'' We note that
\texttt{check\_path\_traversal} is \emph{not} one of the ten mandatory base
sanitizers (Section~\ref{sec:attribution}); its selection is attributable to
the LLM reasoning over the tool's documented file-write semantics.

\para{What \project{} generates ($G$).}
Because the selected sanitizers are file-oriented, the adversarial-generation
stage targets the path-bearing argument directly. Among the ten seeds the LLM
produced for this tool, one places a directory-traversal string into
\texttt{outputPath} (recorded verbatim by the fuzzer):

\begin{lstlisting}[basicstyle=\ttfamily\small,breaklines=true,columns=fullflexible,linewidth=\linewidth]
{ "filePath": "<xlsx>", "sheetName": "Sheet1",
  "outputPath": "../../../../tmp/evil.json",
  "format": "array", "dateFormat": "..." }
\end{lstlisting}

\noindent The traversal payload is generated by the model from the tool's
semantics, not drawn from a static wordlist: the filename is the model's own.
Coverage-guided mutation then explores around this seed (e.g.\ producing
truncated variants such as \texttt{../../../../tmp/evil.jso}), but the seed
already expresses the semantic violation the oracle needs.

\para{The sink.}
The implementation validates only the \emph{type} of \texttt{outputPath} and
then writes to it unconditionally:

\begin{lstlisting}[basicstyle=\ttfamily\small,breaklines=true,columns=fullflexible,linewidth=\linewidth]
// validation: type only
(args.outputPath === undefined || typeof args.outputPath === 'string')
// sink (export-to-json.ts)
if (outputPath) {
  writeFileSync(outputPath, jsonContent, 'utf-8');
}
\end{lstlisting}

\noindent The value flows unmodified from the LLM-supplied argument into
\texttt{writeFileSync}. There is no canonicalisation, no containment check
against a working directory, and no shell---so a conventional
command-injection or shell-metacharacter check finds nothing to flag.

\para{How \project{} detects the violation ($O$).}
Detection here is by \emph{outcome divergence}, not by intercepting the write.
\texttt{excel-reader-mcp} is a Node.js server, and \project{}'s runtime
introspection tracers (Table~\ref{tab:sanitizers}) instrument the Python
runner; on a JS target no tracer observes the \texttt{writeFileSync} call.
Instead, \project{} recognises the danger in the \emph{argument} before
execution---\texttt{check\_path\_traversal} flags the traversal string as
adversarial---and then observes \emph{after} execution that the tool returned a
success response with no exception. Its semantic oracle
(\texttt{on\_trigger:~raise\_exception}) expects a tool to reject an input its
selected sanitizer deemed adversarial; a silent success is the contract
violation, reported as: ``Sanitizers \texttt{[check\_path\_traversal]} flagged
this input as malicious, but the target tool executed successfully without
raising any exception.'' This is the general shape of a \emph{silent} semantic
bug: the failure is not that the tool crashed, but that it accepted and acted
on an argument it should have refused.

\para{Validated consequence.}
Because \texttt{outputPath} is under caller control and unconstrained, the
write escapes the tool's intended output location. We confirmed the capability
by directed re-testing against the server over stdio: an absolute path outside
the working directory is written, a relative \texttt{../} path escapes the
working directory, and---supplying a spreadsheet whose cells carry the desired
bytes---an \emph{existing} file is overwritten with attacker-influenced content
(the serialised cell values). The consequence is thus external control of a
file path leading to arbitrary file write and overwrite
(CWE-73; the \texttt{../} escape is CWE-22), where both the destination and the
written content are influenced by the untrusted, LLM-supplied argument.
\TODO{Ananto: we validated write, relative escape, and overwrite-with-content.
We did NOT validate parent-directory creation or storage/inode exhaustion; do
not add a DoS (CWE-400) claim unless separately measured (Pierre's caution).}

\para{Why this matters for the trust boundary.}
The argument is type-correct, the invocation avoids a shell, and the call
succeeds---so every signal a traditional harness relies on says the tool is
behaving. The bug exists only relative to the tool's \emph{documented meaning}:
a tool that exports a sheet to a JSON file should not be a primitive for
overwriting arbitrary files on the host. \project{} reaches that judgement from
the tool's own metadata---selecting a file-path sanitizer, generating a
path-bearing adversarial input, and flagging the silent acceptance---which is
exactly the argument that LLM-generated tool arguments must be treated as
untrusted at the tool boundary.
\TODO{Ananto: disclose to the maintainer before this names the repo in a
camera-ready (sec10\_ethical). Repo is public, 19 stars, commit 2c41e8a9.}
